\documentclass[12pt,a4paper,openany]{book}

\usepackage{lmodern}
\usepackage{xcolor}
\usepackage{xcolor}
\definecolor{vert1}{rgb}{0.0,0.3.9,0.0}
\definecolor{bleu}{rgb}{0,0,0.5}
\definecolor{bleu3}{rgb}{1,0.2,0.2}
\definecolor{grisgris}{gray}{0.4}
\definecolor{grisclair}{HTML}{E7E7E7}
\definecolor{grisfonce}{HTML}{A5A5A5}
\definecolor{rougeUPS}{rgb}{0.6, 0.3, 0.3}

\fboxsep =0pt \parindent =0pt\parskip =12pt



\usepackage[utf8]{inputenc}
\usepackage[T1]{fontenc}
\usepackage[francais]{babel}
\usepackage[top=1.7cm, bottom=1.7cm, left=1.7cm, right=1.7cm]{geometry}
\usepackage{verbatim}
\usepackage[urlbordercolor={1 1 1}, linkbordercolor={1 1 1}, linkcolor=vert1, urlcolor=bleu, colorlinks=true]{hyperref}
\usepackage{tikz} %Vectoriel
\usepackage{listings}
\usepackage{fancyhdr}
\usepackage{multido}
\usepackage{amssymb}

\newcommand{\titre}{Les réseaux informatiques}
\newcommand{\sigle}{reseau}
\newcommand{\semestre}{3}

\definecolor{gris1}{gray}{0.40}
\definecolor{gris2}{gray}{0.55}
\definecolor{gris3}{gray}{0.65}
\definecolor{gris4}{gray}{0.50}
\definecolor{vert}{rgb}{0,0.4,0}
\definecolor{violet}{rgb}{0.65, 0.2, 0.65}
\definecolor{bleu1}{rgb}{0,0,0.8}
\definecolor{bleu2}{rgb}{0,0.2,0.6}
\definecolor{bleu3}{rgb}{0,0.2,0.2}
\definecolor{rouge}{HTML}{F93928}


\lstdefinelanguage{algo}{%
   morekeywords={%
    %%% couleur 1
		importer, programme, glossaire, fonction, procedure, constante, type, 
	%%% IMPORT & Co.
		si, sinon, alors, fin, tantque, debut, faire, lorsque, fin lorsque, 
		declenche, declencher, enregistrement, tableau, retourne, retourner, =, pour, a,
		/=, <, >, traite,exception, 
	%%% types 
		Entier, Reel, Booleen, Caractere, Réél, Booléen, Caractère,
	%%% types 
		entree, maj, sortie,entrée,
	%%% types 
		et, ou, non,
	},
  sensitive=true,
  morecomment=[l]{--},
  morestring=[b]',
}

\lstset{language=algo,
    %%% BOUCLE, TEST & Co.
      emph={importer, programme, glossaire, fonction, procedure, constante, type},
      emphstyle=\color{bleu2},
    %%% IMPORT & Co.  
	emph={[2]
		si, sinon, alors, fin , tantque, debut, faire, lorsque, fin lorsque, 
		declencher, retourner, et, ou, non,enregistrement, retourner, retourne, 
		tableau, /=, <, =, >, traite,exception, pour, a
	},
      emphstyle=[2]\color{bleu1},
    %%% FONCTIONS NUMERIQUES
      emph={[3]Entier, Reel, Booleen, Caractere, Booléen, Réél, Caractère},
      emphstyle=[3]\color{gris1},
    %%% FONCTIONS NUMERIQUES
      emph={[4]entree, maj, sortie, entrée},	
      emphstyle=[4]\color{gris1},
}
\lstdefinelanguage{wl}{%
   morekeywords={%
    %%% couleur 1
		importer, programme, glossaire, fonction, procedure, constante, type, 
	%%% IMPORT & Co.
		si, sinon, alors, fin, TANTQUE, tantque, FIN, PROCEDURE, debut, faire, lorsque, 
		fin lorsque, declenche, declencher, enregistrement, tableau, retourne, retourner, =, 
		/=, <, >, traite,exception, 
	%%% types 
		Entier, Reel, Booleen, Caractere, Réél, Booléen, Caractère,
	%%% types 
		entree, maj, sortie,entrée,
	%%% types 
		et, ou, non,
	},
  sensitive=true,
  morecomment=[l]{//},
  morestring=[b]',
}

\lstset{language=wl,
    %%% BOUCLE, TEST & Co.
      emph={importer, programme, glossaire, fonction, procedure, constante, type},
      emphstyle=\color{bleu2},
    %%% IMPORT & Co.  
	emph={[2]
		si, sinon, alors, fin , tantque, debut, faire, lorsque, fin lorsque, 
		declencher, retourner, et, ou, non,enregistrement, retourner, retourne, 
		tableau, /=, <, =, >, traite,exception
	},
      emphstyle=[2]\color{bleu1},
    %%% FONCTIONS NUMERIQUES
      emph={[3]Entier, Reel, Booleen, Caractere, Booléen, Réél, Caractère},
      emphstyle=[3]\color{gris1},
    %%% FONCTIONS NUMERIQUES
      emph={[4]entree, maj, sortie, entrée},	
      emphstyle=[4]\color{gris1},
}
\lstdefinelanguage{css}{%
   morekeywords={%
    %%% couleur 1
		background, image, repeat, position, index, color, border, font, 
		size, url, family, style, variant, weight, letter, spacing, line, 
		height, text, decoration, align, indent, transform, shadow, 
		background, image, repeat, position, index, color, border, font, 
		size, url, family, style, variant, weight, letter, spacing, line, 
		height, text, decoration, align, indent, transform, shadow, 
		vertical, align, white, space, word, spacing,attachment, width, 
		max, min, margin, padding, clip, direction, display, overflow,
		visibility, clear, float, top, right, bottom, left, list, type, 
		collapse, side, empty, cells, table, layout, cursor, marks, page, break,
		before, after, inside, orphans, windows, azimuth, after, before, cue, 
		elevation, pause, play, during, pitch, range, richness, spek, header, 
		numeral, punctuation, rate, stress, voice, volume,
	%%% types 
		left, right, bottom, top, none, center, solid, black, blue, red, green,
	},
  sensitive=true,
  sensitive=true,
  morecomment=[s]{/*}{*/},
  morestring=[b]',
}
\lstset{language=css,
    %%% BOUCLE, TEST & Co.
      emph={
		background, image, repeat, position, index, color, border, font, 
		size, url, family, style, variant, weight, letter, spacing, line, 
		height, text, decoration, align, indent, transform, shadow, 
		background, image, repeat, position, index, color, border, font, 
		size, url, family, style, variant, weight, letter, spacing, line, 
		height, text, decoration, align, indent, transform, shadow, 
		vertical, align, white, space, word, spacing,attachment, width, 
		max, min, margin, padding, clip, direction, display, overflow,
		visibility, clear, float, top, right, bottom, left, list, type, 
		collapse, side, empty, cells, table, layout, cursor, marks, page, break,
		before, after, inside, orphans, windows, azimuth, after, before, cue, 
		elevation, pause, play, during, pitch, range, richness, spek, header, 
		numeral, punctuation, rate, stress, voice, volume,
	  },
      emphstyle=\color{bleu2},
    %%% FONCTIONS NUMERIQUES
      emph={[3]
		left, right, bottom, top,none, solid, black, blue, green,
		  },
      emphstyle=[3]\color{bleu3},
    %%% FONCTIONS NUMERIQUES
}

\lstset{language=SQL,
    %%% BOUCLE, TEST & Co.
      emph={INSERT, UPDATE, DELETE, WHERE, SET, GROUP, BY, ORDER, REFERENCES},
      emphstyle=\color{bleu2},
    %%% IMPORT & Co.  
	emph={[2]
		if, end, begin, then, for, each, else, after, of, on, to
	},
      emphstyle=[2]\color{bleu1},
    %%% FONCTIONS NUMERIQUES
      emph={[3]Entier, Reel, Booleen, Caractere, Booléen, Réél, Caractère},
      emphstyle=[3]\color{gris1},
    %%% FONCTIONS NUMERIQUES
      emph={[4]entree, maj, sortie, entrée},	
      emphstyle=[4]\color{gris1},
}
\lstdefinelanguage{ARM}{%
   morekeywords={%
   ADD, SUB, MOV, MUL, RSB,CMP, BLS, BLE, B,BHI,LDR,
   BGE, RSBLT, BGT, BEQ, BNE,BLT,BHS,STR,STRB
	},
  sensitive=true,
  morecomment=[l]{@},
  morestring=[b]',
}

\lstset{ % general style for listings 
   numbers=left 
   , literate={é}{{\'e}}1 {è}{{\`e}}1 {à}{{\`a}}1 {ê}{{\^e}}1 {É}{{\'E}}1 {ô}{{\^o}}1 {€}{{\euro}}1{°}{{$^{\circ}$}}1 {ç}{ {c}}1 {ù}{u}1
	, extendedchars=\true
   , tabsize=2 
   , frame=l
   , framerule=1.1pt
   , linewidth=520px
   , breaklines=true 
   , basicstyle=\footnotesize\ttfamily 
   , numberstyle=\tiny\ttfamily 
   , framexleftmargin=0mm 
   , xleftmargin=0mm 
   , captionpos=b 
	, keywordstyle=\color{bleu2}
	, commentstyle=\color{vert}
	, stringstyle=\color{rouge}
	, showstringspaces=false
	, extendedchars=true
	, mathescape=true
} 
%	\lstlistoflistings
%	\addcontentsline{toc}{part}{List of code examples}
 %prise en charge du langage C 
\date{\today}

\makeindex
\lfoot{Université Toulouse III -- Paul Sabatier}
\rfoot{\sigle\semestre}
%\rfoot{}
\cfoot{}
\makeglossary
\makeatletter
\def\clap#1{\hbox to 0pt{\hss #1\hss}}%
\def\ligne#1{%
\hbox to \hsize{%
\vbox{\centering #1}}}%
\def\haut#1#2#3{%
\hbox to \hsize{%
\rlap{\vtop{\raggedright #1}}%
\hss
\clap{\vtop{\centering #2}}%
\hss
\llap{\vtop{\raggedleft #3}}}}%
\def\bas#1#2#3{%
\hbox to \hsize{%
\rlap{\vbox{\raggedright #1}}%
\hss \clap{\vbox{\centering #2}}%
\hss
\llap{\vbox{\raggedleft #3}}}}%
\def\maketitle{%
\thispagestyle{empty}\vbox to \vsize{%
\haut{}{\@blurb}{}

\vfill
\vspace{1cm}
\begin{flushleft}
\usefont{OT1}{ptm}{m}{n}
\huge \@title
\end{flushleft}
\par
\hrule height 4pt
\par
\begin{flushright}
\usefont{OT1}{phv}{m}{n}
\Large \@author
\par
\end{flushright}
\vspace{1cm}
\vfill
\vfill
\bas{}{\@location, le \@date}{}
}%
\cleardoublepage
}
\def\date#1{\def\@date{#1}}
\def\author#1{\def\@author{#1}}
\def\title#1{\def\@title{#1}}
\def\location#1{\def\@location{#1}}
\def\blurb#1{\def\@blurb{#1}}
\date{\today}
\author{}
\title{}
\location{Amiens}\blurb{}
\makeatother
\title{\titre}
\author{Semestre \semestre}

\location{Toulouse}
\blurb{%
Université Toulouse III -- Paul sabatier\\
L2 Informatique\\
}%



%\title{Cours \\ \titre}
%\date{\today\\ Semestre \semestre}

%\lhead{Cours: \titre}
%\chead{}
%\rhead{\thepage}

%\lfoot{Université Paul Sabatier Toulouse III}
%\cfoot{\thepage}
%\rfoot{\sigle\semestre}

\pagestyle{fancy}
\renewcommand{\chaptermark}[1]{\markboth{\bsc{\chaptername~\thechapter{} :} #1}{}}
\renewcommand{\sectionmark}[1]{\markright{\thesection{ #1}}}
\renewcommand{\headrulewidth}{0.3pt}
\renewcommand{\footrulewidth}{0.3pt}

\fancyhf{}
\fancyhead[LE]{\leftmark}
\fancyhead[RO]{\rightmark}
\fancyfoot[LE,RO]{--~\thepage~--}
\fancyfoot[LO]{\titre{}~~---~~\sigle{}\semestre{}}
\fancyfoot[RE]{Antoine de \bsc{Roquemaurel}}

%% Cas des premières pages de chapitre
\fancypagestyle{plain}{%
	\fancyhf{}%
	\fancyfoot[L]{\titre{}~~---~~\sigle{}\semestre{}}
	\fancyfoot[R]{--~\thepage~--}
	\renewcommand{\headrulewidth}{0pt}
	\renewcommand{\footrulewidth}{0.3pt}
}
\makeatletter
\renewcommand*{\lstlistlistingname}{Liste des codes sources}
\renewcommand\listoffigures{%
    \chapter{\listfigurename}%
      \@mkboth{\MakeUppercase\listfigurename}%
              {\MakeUppercase\listfigurename}%
       \@starttoc{lof}%
    }
    \renewcommand\listoftables{%
    \chapter{\listtablename}%
    \@mkboth{\MakeUppercase{\listtablename}}%
            {\MakeUppercase{\listtablename}}%
    \@starttoc{lot}
    }

    \renewcommand\lstlistoflistings{%
    \begingroup
    \chapter{\lstlistlistingname}%
    \parskip\z@\parindent\z@\parfillskip \z@ \@plus 1fil%
    \@starttoc{lol}%
    \endgroup
    }
	\makeatother

\newcommand{\remarque}[1]{
	\begin{center}
	\medskip
	\colorbox{remarque}{
		\begin{minipage}{0.85\textwidth}\medskip\includegraphics[height=10px]{images/remarque.png} #1 \medskip\end{minipage}
	}
	\medskip
	\end{center}
}

\newcounter{exemples}

\newenvironment{exemple}[1]{
   \vspace{-2mm}

\refstepcounter{exemples}
   \begin{center}
	\medskip
      \begin{minipage}{0.9\linewidth}
}{%
~
      \end{minipage}
   \end{center}~
   \vspace{-2mm}
}%

\newcommand{\captionExemple}[1]{
	\begin{center}{\bsc{Exemple} \thechapter.\arabic{exemples}~--~}#1\end{center}
}

\DeclareTextFontCommand{\policeGlossaire}{\fontfamily{lmss}\selectfont}
\DeclareTextFontCommand{\policePackage}{\fontfamily{phv}\selectfont}
\DeclareTextFontCommand{\policeTitre}{\fontfamily{ptm}\selectfont}
\newcommand{\policeCode}[1]{\texttt{#1}}

\newcommand{\sectionfont}{%
	\fontencoding{\encodingdefault}%
	\fontfamily{pag}%
	\fontseries{bc}%
	\fontshape{n}%
	\selectfont
}

% numéro du chapitre
\DeclareFixedFont{\chapnumfont}{T1}{phv}{b}{n}{80pt}
% pour le mot « Chapitre »
\DeclareFixedFont{\chapchapfont}{T1}{phv}{b}{n}{16pt}
% pour le titre
\DeclareFixedFont{\chaptitfont}{T1}{phv}{b}{n}{24.88pt}


\makeatletter
\def\thickhrulefill{\leavevmode \leaders \hrule height 1ex \hfill \kern \z@}
%% \chapter
\def\@makechapterhead#1{%
  \reset@font
  \parindent \z@
  \vspace*{10\p@}%
  \hbox{%
    \vbox{%
      \advance\hsize by -2cm
      \hrule height 0.4pt depth 0pt width \hsize
      \par
      \vskip 6pt%
      \hspace{20pt}%
      \parbox{420pt}{%
        \LARGE \bfseries #1
		}%
      \par
      \vskip 6pt%
      \hspace{20pt}%
      \hrule height 0.4pt depth 0pt width \hsize
	  \vspace{-30pt}
      }%
    \vbox{%
      \hsize=1.5cm%
      \begin{tabular}{c}
        \scshape \large \strut \@chapapp{} \\
        \colorbox{black}{\vbox{\hbox{\vbox to 1mm{}}\hbox{
			\color{white} \LARGE \bfseries \hspace{1mm}\thechapter\hspace{1mm}
		}\hbox{\vbox to 2cm{}}}}%
      \end{tabular}%
      }%
    }%
  \vskip 20\p@
}
%% \chapter*
\def\@makeschapterhead#1{%
  \reset@font
  \parindent \z@
  \vspace*{10\p@}%
  \hbox{%
    \vbox{%
      \advance\hsize by -0cm
      \hrule height 0.4pt depth 0pt width \hsize
      \par
      \vskip 6pt%
      \hspace{20pt}%
      \parbox{420pt}{%
        \LARGE \bfseries #1
		}%
      \par
      \vskip 6pt%
      \hspace{20pt}%
      \hrule height 0.4pt depth 0pt width \hsize
      }%
    }%
  \vskip 20\p@

}

\newlength{\sectiontitleindent}
\newlength{\subsectiontitleindent}
\newlength{\subsubsectiontitleindent}
\setlength{\sectiontitleindent}{-1cm}
\setlength{\subsectiontitleindent}{-.5cm}
\setlength{\subsubsectiontitleindent}{-.25cm}

\renewcommand{\section}{%
	\@startsection%
	{section}%
	{1}%
	{\sectiontitleindent}%
	{-3.5ex plus -1ex minus -.2ex}%
	{2.3ex plus.2ex}%
	{\sectionfont\Large}
}
\renewcommand{\subsection}{%
	\@startsection%
	{subsection}%
	{2}%
	{\subsectiontitleindent}%
	{-3.5ex plus -1ex minus -.2ex}%
	{2.3ex plus.2ex}%
	{\sectionfont\large}
}

\renewcommand{\subsubsection}{%
	\@startsection%
	{subsubsection}%
	{3}%
	{\subsubsectiontitleindent}%
	{-3.5ex plus -1ex minus -.2ex}%
	{2.3ex plus.2ex}%
	{\sectionfont\normalsize}
}

\makeatother



\begin{document}
	\setcounter{tocdepth}{2}
	\setcounter{secnumdepth}{3}
	\maketitle
	\tableofcontents
	\part{Introduction aux réseaux informatiques}
	\chapter{Histoire de la communication}
		Les humains ont toujours voulu \textbf{communiquer plus vite et plus loin}, ceci en utilisant des codes, alphabets, langages, \ldots
		
		\exemple{Les Gaulois, écrit Jules César dans ``La guerre des Gaules'', avec la voix de champ en champ pouvaient transmettre une nouvelle à 240km de distance en une journée.

			Les Grecs, en utilisant des flambeaux disposés de façon à indiquer les lettres de l'alphabet communiquaient, au temps d'Alexandre, de l'Inde à la Grèce en 5 jours.
			}	

			Le concept de la communication n'a pas changé de nos jours, nous avons toujours un système de codage afin que l'émetteur et le destinataire puisse communiquer. Cependant les supports de la communication ont changé afin de gagner en rapidité (ondes radio, fibre optique\ldots)\footnote{Les supports de communication présentent tous des caractéristiques techniques.}

			\begin{description}
				\item[1464] Poste Royale (Louis XI)\\
					L'inconvénient principal était le temps de transmission.
				\item[1794] Télégraphe optique (Chappe)\\
			Les inconvénients du télégraphe optique sont la visibilité et l'atténuation\ldots Cependant, nous procédons de la même façon, nous utilisons un système de relais: c'est un fondamental.
				\item[1832] Télégraphe \'Electrique (Shilling)
				\item[1837] Code Télégraphique (Morse) et création de l'administration du Télégraphe
				\item[1854] 1$^{er}$ projet de téléphone (Bourseul)
				\item[1860] Lois de l'électromagnétisme (Maxwell)
				\item[1876] Brevet du Téléphone (Bell)
				\item[1887] \'Etude sur les ondes Radioélectriques (Hertz)
				\item[1889] Nationalisation de la société Français de Téléphone
				\item[1892] \'Etude sur la Radiodiffusion (Crooker)
				\item[1896] Liaison de TSF (Marconi)
				\item[1897] \'Emission Radio au Panthéon
				\item[1901] Monopole d'état sur la radiodiffusion
				\item[1915] Téléphone automatique
				\item[1917] Télégraphe de Baudot
				\item[1943] Premier calculateur électronique. Début de l'ère du traitement électronique de l'information: \textbf{Informatique}, 
					suivit de la volonté d'obtenir un moyen de télécommunications entre les équipements Informatique: \textbf{Réseaux Informatiques}.\\ Ainsi, un support de communication, nécessitait un réseau différent (Son $\Rightarrow$ Radio, Image $\Rightarrow$ TV, Texte $\Rightarrow$ Télégraphe, \ldots). Une fusion va se produire.
			\end{description}

			Composants $\Rightarrow$ Signal $\Rightarrow$ Équipements $\Rightarrow$ Protocoles $\Rightarrow$ Architectures $\Rightarrow$ Services.

	\chapter{\'Evolution des réseaux}
			De la même manière que la téléphonie et le télégraphe, nous sommes passé d'une phase expérimentale à une phase d'utilisation. Ainsi l'Informatique à beaucoup évolué. Cette évolution à été progressive, il y a eu plusieurs étapes qui ont marqués les réseaux de communication.
			\paragraph{Coûts des équipements Informatiques / Coûts de la Communication} À l'origine seul les grands comptes étaient capable d'avoir des équipements informatiques. Ainsi les SSI\footnote{Société de Service en Informatiques} sont nées.
			\paragraph{Système de Télétraitement} Ces systèmes ont été destiné aux entreprise, afin qu'a distance elles puissent utiliser la puissance d'un calculateur qui était géographiquement loin. Une première structure de réseau Informatique fut créée.
			\remarque{Nous sommes en train de revenir à cette solution créée 40 ans auparavant: Le cloud computing}

			\section{Les équipements créés}
			Afin de construire ces structures de réseaux de communication nous avons mis en place des équipements :
			\begin{description}
				\item[Processeur Frontal de Communication\footnote{FEP: Front End Processor}]
				\item[Multiplexeurs et concentrateurs] Équipement de partage du support de communication, permettent d'avoir des nœuds de communication.
				\item[Liaisons Spécialisées] Nous avions besoin d'un réseau spécialisé afin d'interconnecter les appareils, pour les connections point à point.
				\item[Modem] Pour les trafics de grande ligne, il fut choisir d'utiliser un réseau déjà existant, le téléphone.
			Cependant, le signal à transmettre doit être adapté au support de transmission, on va donc utiliser un adaptateur qui permettra de faire 
			passer le signal sur le réseau téléphonique : le modem.
				\item[Commutateurs] Pour avoir une connexion la plus rapide possible, nous avions besoin d'un algorithme de routage afin de passer par un chemin en fonction du trafic présent sur la ligne: le routeur.
				\item[Protocole de communication] Permet de faire dialoguer deux machines entre elles, elle doivent utiliser le même protocole afin de se comprendre syntaxiquement et sémantiquement.
			\end{description}
			
			\section{Démocratisation de l'Informatique}
			\begin{description}
				\item[1970] La genèse des protocoles de communication date des années 1970. En réseau, rien n'a été inventé de nouveau, cela à surtout été des progrès technologiques : rapidité, miniaturisation, coûts et donc démocratisation. Les premiers mini-calculateurs.
				\item[1980] Début de l'informatique personnelle et mise en \oe{}uvre des réseaux locaux.
				\item[1990] Applications de l'Internet, premiers mobiles et satellites. 
			\end{description}

	\chapter{Classification}
		Le critère de classification est la distance entre les entités communicantes.
		\begin{enumerate}
			\item Architectures des calculateurs / Architecture de communication
			\item LAN\footnote{Local Area Network} ou RLE\footnote{Réseau Local d'Entreprise}
			\item WAN\footnote{Wide Area Network} ou RLD\footnote{Réseau Longue Distance}
			\item DAN\footnote{Departmental Area Network}, MAN\footnote{Metropolitain Area Network}\ldots
		\end{enumerate}
		\begin{tabular}{| c c c |}
			\hline
			Distance & & Exemple\\
			1m & Square meter & PAN\\
			\hline
			10m & Room & LAN\\
			\hline
			100m & Building & LAN\\
			\hline
			1km & Campus & LAN, DAN\\
			\hline
			100 km & City & WAN\\
			\hline
			1000km & Continent & WAN\\
			\hline
			10 000km & Planet & WAN, The Internet\\
			\hline
		\end{tabular}

		Les autres critères de classifications : 
		\begin{itemize}
			\item Débit
			\item Architecture (OSI,TCP/IP)
			\item \ldots
		\end{itemize}

		Les critères de classifications pour un LAN: 
		\begin{itemize}
			\item PABX
			\item Bureautique
			\item Industriel
			\item Large bande
			\item \ldots
		\end{itemize}
	\chapter{Topologie}
	\chapter{Normalisation}
	\chapter{Architecture de communication}
		\begin{description}
			\item[Transmission physique] Correspond au supports, au type d'encodage, aux liaisons, toute l'architecture physique.
			\item[Contrôle d'erreurs]  Vérifier que les paquets sont bien arrivés.
			\item[Contrôle de flux]  S'assurer que l'émetteur n'aille pas trop vite par rapport au récepteur
			\item[Routage]  en cas de nœud de communication, choisir le chemin le plus rapide ceci en fonction du trafic. 
			\item[Régulation de flux (congestion)] Réguler le flux pour éviter la congestion\footnote{Peut être assimiler aux bouchons}, il préfère la prévention afin d'éviter les bouchons. 
			\item[Séquencement]  Les fichiers sont découpés en plusieurs paquets, en effet un paquet à une taille maximum, le séquencement réassemble les paquets afin de reconstituer le fichier grâce à une numérotation des paquets..
			\item[Contrôle de bout en bout] Vérifier que le fichier à bien été reconstitué. 
			\item[Gestion du dialogue] 
			\item[Reprise sur incidents]Cela permet aussi de gérer un arrêt de la connexion réseau afin de reprendre le transfert à l'endroit où il s'était arrêté
			\item[Transformation de l'information] Codage de l'information (avec le code ASCII par exemple), compression (codecs), sécurité de l'information (cryptage) 
			\item[Synchronisation des processus] Sémantique de l'application, c'est-à-dire quelle opération au niveau applicatifs (Renommer un fichier, créer un répertoire, \ldots)
		\end{description}

		Des \textbf{architectures normalisées} ont été mis en place par les opérateurs de Télécommunications (X21, X25, ISDN,\ldots).\\
		Des \textbf{architectures propriétaires} ont également être mise en place par les constructeurs informatiques (SNA, DNA, DSA).\\
		\begin{description}
			\item[1977] \textbf{ISO} constitue un comité pour la normalisation dans le domaine des télécommunications et de l'interconnexion des systèmes.\\
			\item[1984] ISO 7948 référence CCITT X.200 (ITU)
			\item[OSI] Cadre fonctionnel -- Le modèle de référence. Les objectifs du modèle OSI sont les suivants : 
				\begin{itemize}
					\item Décomposer (décomposition fonctionnelle)
					\item Structurer
					\item Assurer l'indépendance vis à vis du matériel et du logiciel.
				\end{itemize}
		\end{description}
	\chapter{Modèle OSI}
	\section{Les couches}
\begin{verbatim}
      +-------------------+ Exemple  
7 N+4 |   Application     | HTTP   
      +-------------------+     
6 N+3 |   Présentation    | HTML, MPEG     
      +-------------------+              SERVICES APPLICATIFS (Couches hautes)
5 N+2 |   Session         |   
      +-------------------+ ===========================================================
4 N+1 |   Transport       | TCP 
      +-------------------+    
3 N   |   Réseau          | IP    
      +-------------------+                 SERVICES TRANSPORTS(Couches basses)
2 N-1 |   Liaison         | Wifi    
      +-------------------+    
1 N-2 |   Physique        | RJ45  
      +-------------------+  
\end{verbatim}
\begin{verbatim}
	Système A                                               Système B
+--------------+                                           +---------------+
| Layer (N+1)  | <---------------------------------------> |  Layer (N-1)  |
+--------------+                                           +---------------+       
      ^                                                         ^
      |                                                         | 
  Service Interface                                       Sercice Interface
    =======                                                   =======
      | SAP                                                     | SAP
      |                                                         | 
      v                                                         v
+---------------+	                                        +---------------+	   
|  Layer (N)    |  <--------------------------------------> |   Layer (N)   |
+---------------+	                                        +---------------+
\end{verbatim}
\newpage
\section{Communication entre systèmes}
\begin{verbatim}
  +-------------------+                                     +-------------------+
7 |   Application     |   <------------------------------>  |  Application      | 7
  +-------------------+                                     +-------------------+
6 |   Présentation    |   <------------------------------>  |  Présentation     | 6 
  +-------------------+                                     +-------------------+
5 |   Session         |   <------------------------------>  |  Session          | 5 
  +-------------------+                                     +-------------------+
4 |   Transport       |   <------------------------------>  |  Transport        | 4
  +-------------------+                                     +-------------------+
3 |   Réseau          |   <------------------------------>  |  Réseau           | 3
  +-------------------+                                     +-------------------+ 
2 |   Liaison         |   <------------------------------>  |  Liaison          | 2
  +-------------------+                                     +-------------------+
1 |   Physique        |   <------------------------------>  |  Physique         | 1
  +-------------------+                                     +-------------------+

\end{verbatim}

\section{4 primitives de Service} 
\begin{description}
	\item[Requête] une entité sollicite un service pour faire une activité.
	\item[Indication] Informe d'un événement
	\item[Réponse] Réponse à l'événement.
	\item[Confirmation] Informe de la demande de service
\end{description}
Modes de communication:
\begin{itemize}
	\item Mode connexté
	\item Mode non contecté
\end{itemize}

Trois phases pour le mode connecté
\begin{itemize}
	\item \'Etablissement de la connexion avec négociation entre les deux entités (N+1) et le service (N)
	\item Transfert de données entre entités (N+1) sur la connexion (N) avec séquencement.
	\item Libération de la connexion
\end{itemize}
Multiplexage de connexion et éclatement de connexions
	\chapter{Approche métiers}

	\part{Transmission des données}
	\chapter{Codage de l'Information}
	\begin{itemize}
		\item Communication = Transmission + Compréhension
		\item Langages, Ecrit,\ldots
		\item Télégraphie, Morse, Baudot
		\item Télex, terminal ``Numérique''
		\item Codage Informatique : suite de chiffres binaires (bits ou binary digits) 0 et 1
	\end{itemize}
	\appendix
	\chapter{Exo}
	\subsection{Comment appelle-t-on un réseau à l'échelle d'une entreprise ?}
	\subsection{Quelles sont les différences entre une topologie point à point et une topologie à diffusion ?}
	\subsection{Quel est l'intérêt de lanormalisation ? Citer 4 organismes de normalisation}
	\subsection{Qu'appelle-t-on système réél? Comment un système réel peut devenir ouvert ?}
	\subsection{Que définit le modèle de référence OSI ?}
	\subsection{Quelle est la signification de la communication horizontale dans le modèle OSI ?}
	\subsection{Dans quelle mesure peut-on parler de communication verticale dans le modèle OSI ?}
	\subsection{Qui du modèle OSI ou de l'architecture TCP/IP est le plus récent ?}
	\subsection{À quel niveau se situe le déplacement proprement dit des données ?}
	\subsection{Quel type d'identifiant est mis en \oe{}uvre au niveau 1 OSI ?}
\end{document}



